\begin{lemma}[Claim 4.3(b2), tree-edge clause]
In the setting of Algorithm (A), suppose that the initial coloring
\(\sigma_n\) is proper, leaves exactly \(e_0\) uncolored, and is stable
on the terminal ladder tree relative to its terminal stored colors and
displayed coloring.  If the restoration trace satisfies
\noderef{ExtensionBaseRestorationInvariant}, then its output
\(\varrho_1\) satisfies, for every exact displayed edge identity \(f\)
of \(T_j\),
\[
             \varrho_1(f)=\sigma_n(f)=\varphi_{j-1}(f).
\]
\end{lemma}
\begin{proof}
Fix an exact displayed edge identity
\(f=e_a\) of \(T_j\), and fix a restoration step \(b<\ell\).  Write
\(t=t_b\), \(v_t=v_{t_b}\), \(\sigma^b\) for the coloring before this
step, and
\[
 C_b=P_{v_t}(\gamma_t,\delta_t,\sigma^b)
\]
for its rooted two-color component.  The initial-segment clause of
\noderef{ExtensionBaseRestorationInvariant} says that \(T_j\) is an
ordered initial segment of the current tree \(T_t\).  Unpacking
\noderef{TreeSequenceInitialSegment}, there is a length inequality and
the vertex and edge entries of \(T_j\) agree entry by entry with the
corresponding entries of \(T_t\).  In particular, the occurrence of
\(e_a\) in \(T_t\) is the same exact edge identity \(f\), not merely an
edge with the same ends.

By the attachment clause of \noderef{TreeSequence}, for some earlier
index \(i<a\) one has
\[
             \operatorname{ends}(f)=\{y_a,y_i\}.              \tag{1}
\]
The entrywise vertex equalities just obtained put the corresponding
occurrences of both \(y_a\) and \(y_i\) among the vertices of \(T_t\),
so \noderef{TreeSequenceVertexSet} gives
\(y_a,y_i\in V(T_t)\).  Equation (1) also says that each of these two
vertices is incident with the exact edge \(f\).

We claim that the edge-node \(\operatorname{inr}(f)\) is not in
\(C_b\).  Suppose to the contrary that it is.  By
\noderef{TwoColorComponent}, this means that the edge-node and the root
vertex-node \(\operatorname{inl}(v_t)\) lie in the same connected
component of the graph in \noderef{TwoColorIncidenceGraph}.  These two
nodes have different sum tags, so a component path between them is
nontrivial.  At its edge-node end it has an incidence adjacency.  The
definition of the two-color incidence graph therefore implies
\[
          \sigma^b(f)=\gamma_t\quad\hbox{or}\quad
          \sigma^b(f)=\delta_t.                              \tag{2}
\]
For each \(x\in\{y_a,y_i\}\), incidence of \(x\) with \(f\), together
with (2), makes \(\operatorname{inl}(x)\) adjacent to
\(\operatorname{inr}(f)\) in that same incidence graph.  Hence both
endpoint vertex-nodes belong to \(C_b\).  The exact-intersection clause
of \noderef{ExtensionBaseRestorationInvariant}, applied separately to
the already established memberships \(y_a,y_i\in V(T_t)\), now gives
\[
                         y_a=v_t=y_i.                         \tag{3}
\]
Substitution of (3) into (1) yields
\(\operatorname{ends}(f)=\{v_t,v_t\}\), contradicting the loopless
law in \noderef{LooplessMultigraph}.  This proves the claim for the
arbitrary step \(b\).

The transition clause of \noderef{ExtensionBaseRestorationTrace} says
that \(\sigma^{b+1}\) is obtained from \(\sigma^b\) by the component
switch rooted at \(v_t\).  Since the edge-node of \(f\) is outside
that component, the outside-component clause in the definition
\noderef{TwoColorComponentSwap} gives
\[
                         \sigma^{b+1}(f)=\sigma^b(f).          \tag{4}
\]
We now induct explicitly on \(m\), for \(0\leq m\leq\ell\), to prove
\(\sigma^m(f)=\sigma_n(f)\).  At \(m=0\), this is the
\(\sigma^0=\sigma_n\) initial-coloring clause of
\noderef{ExtensionBaseRestorationTrace}.  If it holds at
\(m=b<\ell\), equation (4) gives
\(\sigma^{b+1}(f)=\sigma^b(f)=\sigma_n(f)\).  Finite induction reaches
\(m=\ell\); since the output is \(\varrho_1=\sigma^\ell\), it gives
\[
                         \varrho_1(f)=\sigma_n(f).              \tag{5}
\]

Apply \noderef{ExtensionBaseClaimTreeEdgeColors} to the stated indices,
properness, exact-one-uncolored property, and terminal stability of
\(\sigma_n\).  Its displayed-edge conclusion, which retains the exact
edge identity \(f\), gives
\[
                         \sigma_n(f)=\varphi_{j-1}(f).          \tag{6}
\]
Combining (5) and (6) proves the tree-edge part of (b2).
\end{proof}
